% Vista científica exacta materializada desde una rama condicional.
% Fuente: source/demostraciones_primarias/27_04_particula_estadistica.tex; rama: HMTIncludeParticleCore.
% No modifica el contenido matemático de la rama de origen.
\chapter{Sections persistantes et propagation finie}
\label{chap:particula}

La notion mathématique de particule est introduite après la structure de frontière et avant la spécialisation de sa dynamique. Son support primaire est une extension survivante du système d’états TPK\@. Une trajectoire orientée est l’ombre séquentielle de cette extension ; sur elle, les degrés internes forment un fibré discret. Le propagateur, la somme sur les chemins et les représentations de Clifford sont introduits après la construction de la droite d’action. Cette hiérarchie empêche qu’une particule soit définie par le choix externe d’une ligne ou d’un parcours et distingue l’objet mathématique de son identification ultérieure à une espèce physique.

\section{Extension survivante et projection observable}

Soit \((\mathcal S_m,\rho_{m+1,m})\) le système inverse des états qui franchissent les clôtures jusqu’à la profondeur \(m\). Un état global est une famille compatible
\[
 x=(x_m)_m\in\Omega_{\HMT}^{\rm surv}
 :=\varprojlim_m\mathcal S_m.
\]
Une carte séquentielle de profondeur \(R\) produit un parcours
\[
 \gamma_x^{[R]}=\operatorname{sh}_R(x)
 =(c_0\longrightarrow\cdots\longrightarrow c_R).
\]
Le symbole \(\operatorname{sh}\) désigne une lecture d’ombre : différentes cartes peuvent présenter un même état global, pourvu qu’elles conservent le feuillet, l’orientation, le bord et l’enregistrement. Dans le \cref{chap:extension-ruta-caracter} est construit l’enregistrement enrichi qui fait descendre le caractère d’action de ces présentations à l’extension.

\section{Particule comme section compatible d’un fibré discret}

Soit \(\Gamma_{\rm TPK}\) le graphe dont les sommets sont les états complets du noyau de phase topologique et dont les arêtes sont les transformations admissibles. Pour \(x\in\Omega_{\HMT}^{\rm surv}\) et sa trajectoire observée
\[
 \gamma_x=(c_0\longrightarrow c_1\longrightarrow\cdots
 \longrightarrow c_R),
\]
soit \(\operatorname{Typ}(c_t)\) l’ensemble fini des types arithmétiques, des orientations et des mémoires sur l’état \(c_t\), et soit \(E_{c_t}=\mathbb C[\operatorname{Typ}(c_t)]\) l’espace vectoriel complexe libre engendré par ces types. Le fibré discret restreint à \(\gamma_x\) est
\[
 \mathcal B_{\gamma_x}=\bigsqcup_{t=0}^{R}E_{c_t}
 \longrightarrow\{0,\ldots,R\}.
\]
Chaque arête \(c_t\to c_{t+1}\) induit une application linéaire de transport \(\tau_t:E_{c_t}\to E_{c_{t+1}}\).

\begin{definition}[Particule mathématique]
Une particule de Mécanique Dimensionnelle est un tuple
\[
 \mathfrak p=(x,\gamma_x,\mathcal B_{\gamma_x},s,n,U),
\]
où \(x\in\Omega_{\HMT}^{\rm surv}\) est l’extension survivante et \(\gamma_x=\operatorname{sh}_R(x)\) son ombre dans la carte déclarée ; de plus,
\[
 s:\{0,\ldots,R\}\longrightarrow\mathcal B_{\gamma_x},\qquad
 s(t)\in E_{c_t},\qquad s(t+1)=\tau_t(s(t)),
\]
est une section compatible ;
\[
 n:\bigsqcup_{t=0}^{R}\operatorname{Modes}(E_{c_t})
 \longrightarrow\mathbb N
\]
est sa fonction d’occupation, où \(\operatorname{Modes}(E_{c_t}):=\operatorname{Typ}(c_t)\) désigne la base distinguée de l’espace vectoriel libre ; et
\[
 U_{t_2,t_1}:E_{c_{t_1}}\longrightarrow E_{c_{t_2}}
\]
est une famille de propagateurs linéaires qui entrelace les transports, \(U_{t,t}=\operatorname{id}_{E_{c_t}}\), \(U_{t+1,t}=\tau_t\), et vérifie
\[
 U_{t_3,t_2}\circ U_{t_2,t_1}=U_{t_3,t_1}
 \qquad(0\leq t_1\leq t_2\leq t_3\leq R).
\]
Chaque composition est définie entre les fibres correspondantes. L’identification de \(\mathfrak p\) à une espèce physique requiert en outre une représentation du groupe de symétrie pertinent et une application des états HMT vers les états de cette représentation.
\end{definition}

L’extension conserve la généalogie et l’enregistrement ; le parcours en expose la lecture ordonnée ; les degrés internes appartiennent aux fibres, l’occupation à \(n\) et l’évolution à \(U\). La particule est la compatibilité de ces niveaux, et non l’un d’eux pris séparément. En particulier, le parcours n’est pas choisi par la particule : il est une projection de l’extension qui la réalise.
